\documentclass[10pt,a4paper]{amsart}
\usepackage[margin=19mm]{geometry}
\usepackage{amsmath,amssymb,mathtools}
\usepackage[scr=rsfs,cal=euler]{mathalfa}
\usepackage{xfrac}
\usepackage[unicode,hidelinks,bookmarks=false]{hyperref}
\newcommand{\ie}{i.e.\ }
\newcommand{\cf}{cf.\ }
\newcommand{\resp}{respectively\ }
\newcommand{\Subsection}{\subsection}
\providecommand{\nobreakdash}{\nobreak\hbox{-}}
\setlength{\emergencystretch}{2em}
\multlinegap0pt
\usepackage{bm}
\usepackage{bbm}
\usepackage{dsfont}
\usepackage{xspace}
\usepackage{cancel}
\usepackage{mathtools}
\usepackage{amsthm}
\makeatletter
\@ifundefined{theorem}{\newtheorem{theorem}{Theorem}}{}
\@ifundefined{lemma}{\newtheorem{lemma}[theorem]{Lemma}}{}
\makeatother
\makeatletter
\expandafter\ifx\csname proof\endcsname\relax
\newenvironment{proof}[1][Proof]{\noindent\textbf{#1.} }{\ \rule{0.5em}{0.5em}}
\fi
\makeatother
\title[An inequality of Harish-Chandra]{An inequality of Harish-Chandra}
\author{Nolan R Wallach}
\date{Source: arXiv:2505.14960v1 (CC BY 4.0)}
\begin{document}
\maketitle

\noindent\emph{Source and licence.} An inequality of Harish-Chandra. Version arXiv:2505.14960v1 (CC BY 4.0), distributed under the Creative Commons Attribution 4.0 International licence, as recorded in the arXiv licence metadata for this exact version and shown on the abstract page https://arxiv.org/abs/2505.14960v1. Full rights notice: licenses/arxiv-2505.14960-CC-BY-4.0.txt.\par\smallskip

\noindent\textit{Editorial scope.} Counted unit: the Lemma whose source label is \texttt{MyInequality} in the article (source0.tex lines 239--267, source label \texttt{MyInequality}), delivered together with the article's own complete original proof of it (source0.tex lines 269--414, one \texttt{proof} environment of 146 lines). No mathematics is written, repaired or re-derived: statement and proof are the article's own text line for line. Required context retained uncounted: trivial (lines 418--421). Every definition, display and label of those lines is retained unchanged. Omitted from this package and not redistributed: 1--238, 268--268, 415--417, 422--562. Nothing inside a retained excerpt is omitted or altered: every retained body is delivered exactly as the article's lines are written, so no omission receipt of any kind accompanies this package. Citations: the retained excerpts carry no citation command at all, so no bibliography entry of the article is transcribed and no reference is redistributed. Reference adaptation: none was needed and none was made; every reference inside the delivered unit resolves inside it, so no unresolved reference or citation is produced. Printed slips kept literal: no printed word, symbol, hypothesis or number of the counted statement or of its proof was altered. Layout and class: the article's own class is \texttt{article} with packages including amsfonts, amsmath, amssymb, graphicx; the wrapper is the lane's pinned \texttt{amsart} editorial wrapper at 10pt on A4, which restates the theorem-like environments the retained text needs and adds the article's own macro definitions that the retained text needs (none), with extra wrapper packages (bm, bbm, dsfont, xspace, cancel, mathtools). The delivered numbering is the unit's own. Assets: no retained span contains a figure environment, a graphics-inclusion command, a PDF-page inclusion, a table or a TikZ picture; no image file is copied into this package, the package declares no asset dependency and no image file is redistributed with it. Licence: version arXiv:2505.14960v1 of this article is distributed under the Creative Commons Attribution 4.0 International licence, as recorded in the arXiv licence metadata for this exact version and shown on the abstract page https://arxiv.org/abs/2505.14960v1. A full rights notice is delivered at licenses/arxiv-2505.14960-CC-BY-4.0.txt. Characters: every retained excerpt is pure ASCII, so the delivered engine, XeLaTeX with its default Latin Modern text font, reports no missing character.
\begin{lemma}
\label{trivial}If $a\geq1$ and $b\geq0$ and if $0<m\leq1$ then $a^{m}%
+b\geq\frac{1}{2}(a+b)^{m}$.
\end{lemma}

\begin{lemma}
\label{MyInequality}Let $\left(  W,\left\langle ...,...\right\rangle \right)
$ and $(V,(...,...))$ be non-zero finite dimensional Hilbert spaces over
$\mathbb{R}$ such that $W=\oplus_{j=1}^{q}W_{j}$ and $V=\oplus_{_{j=0}}%
^{p}V_{j}$ orthogonal direct sums. If $w\in W$ (resp. $v\in V)$ set $w_{j}$
(resp. $v_{j}$) equal to its orthogonal projection into $W_{j}$ (resp. $V_{j}%
$). Let $\varphi:W\rightarrow V$ be a map satisfying:

If $j=0$ then $\varphi(w)_{0}$ is non-zero and constant.

If $j=1$ then%
\[
\varphi(w)_{1}=T_{1}w_{1}\text{.}
\]


If $1<j\leq q$ then%
\[
\varphi(w)_{j}=T_{j}w_{j}+u_{j}(w_{1}+...+w_{j-1})
\]
with $T_{j}$ an injective linear map of $W_{j}$ into $V_{j}$ for $j=1,...,q$
and $u_{j}$ is a polynomial of degree $d_{j}$ with values in $V_{j}.$

Then there exists $m>0$ and $C>0$ such that
\[
\left\Vert \varphi(w)\right\Vert \geq C(1+\left\Vert w\right\Vert ^{2})^{m}.
\]

\end{lemma}

\begin{proof}
Let $v_{0}=\varphi(w)_{0}$ (for all $w\in W$). Replacing $\varphi$ with
$\frac{1}{\left\Vert v_{0}\right\Vert }\varphi$ we may assume that $\left\Vert
v_{o}\right\Vert =1$. Define%
\[
\varphi_{r}(w)=\sum_{j=0}^{r}\varphi(w)_{j}.
\]
Then the conditions on $\varphi$ imply that if $r\leq q$ then $\varphi
_{r}(w)=\varphi_{r}(w_{1}+...+w_{r})$. We first prove by induction on $r$ that
there exist $C_{r}>0$ and $1\geq m_{r}>0$ such that
\[
\left\Vert \varphi_{r}(w)\right\Vert ^{2}\geq C_{r}(1+\left\Vert
w_{1}+...+w_{r}\right\Vert ^{2})^{m_{r}}.
\]
If $r=1$ then
\[
\varphi_{1}(w_{1})_{0}+\varphi_{1}(w_{1})_{1}=v_{0}+T_{1}w_{1}.
\]
so%
\[
\left\Vert \varphi_{1}(w_{1})\right\Vert ^{2}=1+\left\Vert T_{1}%
w_{1}\right\Vert ^{2}
\]
since $T_{j}$ is injective \thinspace$\left\Vert T_{j}y_{j}\right\Vert
\geq\lambda_{j}\left\Vert y_{j}\right\Vert $. Let $c_{1}$ be the minimum of
$1$ and $\lambda_{1}^{2}$. Then
\[
\left\Vert \varphi_{1}(w_{1})\right\Vert ^{2}\geq c_{1}(1+\left\Vert
w_{1}\right\Vert ^{2}).
\]
Assume the result for $r-1\geq1$. If $r>q$ then since $r-1\geq q$ we have
\[
\left\Vert \varphi_{r}(w)\right\Vert ^{2}=\left\Vert \varphi_{r-1}%
(w)+\varphi(w)_{r}\right\Vert ^{2}\geq\left\Vert \varphi_{r-1}(w)\right\Vert
^{2}\geq C_{r-1}(1+\left\Vert w_{1}+...+w_{q}\right\Vert ^{2})^{m_{r-1}}
\]%
\[
=C_{r-1}(1+\left\Vert w\right\Vert ^{2})^{m_{r-1}}.
\]
So we may assume that $r\leq q$. Proceeding with the inductive step
\[
\left\Vert \varphi_{r}(w_{1}+...+w_{r})\right\Vert ^{2}=\left\Vert
\varphi_{r-1}(w_{1}+...+w_{r-1})\right\Vert ^{2}+\left\Vert T_{r}w_{r}%
+u_{r}(w_{1}+...+w_{r-1}\right\Vert ^{2}.
\]
Let
\[
S_{r}=\{w_{1}+...+w_{r}|\left\Vert T_{r}w_{r}\right\Vert >2\left\Vert
u_{r}(w_{1},...,w_{r-1})\right\Vert \}.
\]
If
\[
w_{1}+...+w_{r}\in S_{r}
\]
then%
\[
\left\Vert T_{r}w_{r}+u_{r}(w_{1}+...+w_{r-1})\right\Vert \geq\frac{1}%
{2}\left\Vert T_{r}w_{r}\right\Vert \geq\frac{\lambda_{r}}{2}\left\Vert
w_{r}\right\Vert .
\]
Thus%
\[
\left\Vert \varphi_{r}(w_{1}+...+w_{r})\right\Vert ^{2}\geq\left\Vert
\varphi_{r-1}(w_{1}+...+w_{r-1})\right\Vert ^{2}+\frac{\lambda_{r}^{2}}%
{4}\left\Vert w_{r}\right\Vert ^{2}
\]%
\[
\geq C_{r-1}^{2}(1+\left\Vert w_{1}+...+w_{r-1}\right\Vert ^{2})^{m_{r-1}%
}+\frac{\lambda_{r}^{2}}{4}\left\Vert w_{r}\right\Vert ^{2}
\]
using the inductive hypothesis with $1\geq m_{r-1}>0$. Setting $L=\min
\{C_{r-1}^{2},\frac{\lambda_{r}^{2}}{4}\}$ then%
\[
\left\Vert \varphi(w_{1}+...+w_{r})\right\Vert ^{2}\geq L((1+\left\Vert
w_{1}+...+w_{r-1}\right\Vert ^{2})^{m_{r-1}}+\left\Vert w_{r}\right\Vert
^{2}).
\]
In Lemma \ref{trivial} take $a=1+\left\Vert w_{1}+...+w_{r-1}\right\Vert ^{2}
$ and $b=\left\Vert w_{r}\right\Vert ^{2}$ then we have%
\[
\left\Vert \varphi(w_{1}+...+w_{r})\right\Vert ^{2}\geq\frac{L}{2}%
((1+\left\Vert w_{1}+...+w_{r-1}\right\Vert )^{2}+\left\Vert w_{r}\right\Vert
^{2})^{m_{r-1}}.
\]


If
\[
w_{1}+...+w_{r}\notin S_{r}
\]
then%
\[
\left\Vert u_{r}(w_{1},...,w_{r-1})\right\Vert \geq\frac{1}{2}\left\Vert
T_{r}w_{r}\right\Vert \geq\frac{\lambda_{r}}{2}\left\Vert w_{r}\right\Vert .
\]
The inductive hypothesis implies that
\[
\left\Vert \varphi(w_{1}+...+w_{r-1})\right\Vert ^{2}\geq C_{r-1}(1+\left\Vert
w_{1}+...+w_{r-1})\right\Vert ^{2})^{m_{r-1}}
\]
also since $u_{r}(w_{1},...,w_{r-1})$ is a polynomial of degree $d_{r}$
\ there exists $M_{r}$ such that
\[
\left\Vert u_{r}(w_{1},...,w_{r-1})\right\Vert \leq M_{r}(1+\left\Vert
w_{1}+...+w_{r-1}\right\Vert ^{2})^{\frac{d_{r}}{2}}.
\]
Thus
\[
\left\Vert u_{r}(w_{1},...,w_{r-1})\right\Vert \leq\frac{M_{r}}{C_{r-1}%
}\left\Vert \varphi_{r-1}(w_{1}+...+w_{r-1})\right\Vert ^{\frac{d_{r}%
}{2m_{r-1}}}.
\]
This implies that
\[
\left\Vert w_{r}\right\Vert ^{2}+C_{r-1}(1+\left\Vert w_{1}+...+w_{r-1}%
)\right\Vert ^{2})^{m_{r-1}}\leq\left(  \frac{2M_{r}}{\lambda_{r}C_{r-1}%
}\right)  ^{2}\left\Vert \varphi_{r-1}(w_{1}+...+w_{r-1})\right\Vert
^{\frac{d_{r}}{2m_{r-1}}}
\]%
\[
+\left\Vert \varphi_{r-1}(w_{1}+...+w_{r-1})\right\Vert
\]
Let $L=\min\{C_{r-1},1\}$ then
\[
\left\Vert w_{r}\right\Vert ^{2}+C_{r-1}(1+\left\Vert w_{1}+...+w_{r-1}%
)\right\Vert ^{2})^{m_{r-1}}\geq L(\left\Vert w_{r}\right\Vert ^{2}%
+C_{r-1}(1+\left\Vert w_{1}+...+w_{r-1})\right\Vert ^{2})^{m_{r-1}})
\]%
\[
\geq L_{1}(1+\left\Vert w_{1}+...+w_{r}\right\Vert ^{2})^{m_{r-1}}.
\]
Thus
\[
L_{1}(1+\left\Vert w_{1}+...+w_{r}\right\Vert ^{2})^{m_{r-1}}\leq
L_{2}\left\Vert \varphi_{r-1}(w_{1}+...+w_{r-1})\right\Vert ^{\max
\{\frac{d_{r}}{2m_{r-1}},1\}}.
\]
We have shown $w_{1}+....+w_{r}\notin S_{r}$ there exists $L_{3}>0$ such that%
\[
\left\Vert \varphi_{r}(w_{1}+...+w_{r})\right\Vert \geq L_{3}(1+\left\Vert
w_{1}+...+w_{r}\right\Vert ^{2})^{\frac{m_{r-1}}{\max\{\frac{d_{r}}{2m_{r-1}%
},1\}}}.
\]
To complete the inductive step take $C_{r}=\min\{L_{3},C_{r-1}\}$ and
$m_{r}=\min\{m_{r-1,}\frac{m_{r-1}}{\max\{\frac{d_{r}}{2m_{r-1}},1\}}\}$.
\end{proof}

\end{document}
